% BEGIN RECOVERY-STRENGTHEN5 2026-10-08
\section{Recovered phase boundary for the all-input route}
The original RECURCLOSE1, CYCLEGAIN1, STREAMCOMP1 and ORBITBATCH1 results
solve composition and accounting only conditionally.  Their strongest
combined consequence is that strict shrinkage need not hold on every
internal typed edge.  A finite, well-founded completed phase may contain
same-scale edges and may then be flattened into one macro transition.

At every phase entry of encoded length $N$, the macro must decide exactly or
emit ordinary SAT exits of lengths $m_i$ with
\[
 \sum_i m_i\le N^{1+o(1)},\qquad
 \max_i m_i\le N^{h+o(1)}
\]
for fixed $h<1$.  All internal discovery, failed attempts, classification,
setup, nested execution, repeated access, replay and combination must total
$N^{1+o(1)}$; peak simultaneous state including same-scale representations
and ancestors must be $N^{p+o(1)}$ for fixed $p<1$, plus one charged exit
buffer.  Under these premises aggregate recursive closure gives one joint
$N^{1+o(1)}$-time, $N^{\max\{h,p\}+o(1)}$-space machine.

This preserves conditional elimination of the starting exponent $K$ and
allows double-wrapper or typed-cycle candidates.  It does not construct the
required all-input phase.  The ARC route still needs an actual finite joint
transition into the verified FLMV region.  P versus NP remains OPEN.
% END RECOVERY-STRENGTHEN5
